\section{2.2 낙관적 초기화와 UCB: 불확실성을 이용한 탐험}
\begin{frame}{두 가지 철학: ``내부 유도형'' 탐험}
  2.1절의 ε-greedy는 탐험 비율을 \textbf{정해진 확률}
  $\varepsilon$ 로 정하는 ``외부에서 조종하는'' 방식.
  \vskip0.5em
  이 절의 두 전략은 정반대 철학 --- \textbf{탐험을 알고리즘 내부에서
  자동으로 유도}:
  \begin{itemize}
    \item \kb{낙관적 초기화(optimistic initialization)}:
      ``모든 팔을 처음엔 과대평가한다''는 \textbf{시작값 하나로}
      순수 탐욕만으로도 초반 탐험이 자동 생성.
    \item \kb{UCB(Upper Confidence Bound)}: ``아직 덜 당겨본 팔의
      추정치는 그만큼 불확실하다''는 점을 \textbf{매 스텝마다}
      계산에 반영.
  \end{itemize}
  \vskip0.5em
  \kq{왜 ``내부 유도형''이 필요한가? ε-greedy의 탐험은
  \textbf{정보와 무관하게} 발생 --- 이미 ``최악''이라 확신한 팔도
  계속 당김. 이 절의 전략은 탐험을 \textbf{불확실성에 비례해서만}
  씀.}
\end{frame}
